% !TeX root = ../papers/nc-p3.tex
\section{任意の余因子曲線の横断次数と核線束の下界}
\label{sec:sectional-cofactor-budget}
\begin{theorem}\label{thm:sectional-cofactor-budget}
被約四辺有限標準正則窓の整な非恒等反対八次 $D$ で、核グラフの滑らかな解消の
入力束を $H$、核束を $L$ とする。このとき $HL\ge5$ である。
階数二以下の被約曲線の既約成分を $C_i$ とし、その一般横断平面上で
余因子イデアルの極大イデアル位数を $r_i$、平面八次の重複度を $m_i$ とすると
\[
         \sum_i r_i m_i\deg C_i\le38.
\]
特に階数二以下の被約総次数を $d$、その中で一般階数一の成分の被約総次数を $j$ として
$d+2j\le19$。成分の個数にも一般点のスキーム重複度にも仮定しない。
\end{theorem}
\begin{proof}[曲線核の跡と四次線形系]
定理\ref{thm:kernel-curve-scroll}の記号で、核像が曲線 $K$ ならその正規化 $B$ の種数を $g$、
核の四切断空間を $U\subset H^0(B,L_B)$、$e=\deg L_B$、$a=e-g-3\ge0$ とする。
正規化母線束の射影を $\pi$、入力相対二束を $E=\pi_*\mathcal O(H)$ と書く。
一般入力平面は各母線と一点で交わるので $HL=e$ である。
標準跡の包含を大域切断でも使うと、底上に有効因子 $\Theta$ があり
\[
   \deg\Theta=2a,\qquad
   \tau H^0(B,\omega_B\otimes\det E\otimes L_B^{-1})\subset U
\]
となる。ここで $\tau$ は $\Theta$ を切る切断である。
実際 $M\subset\nu_*\pi^* L_B$ を双対にして二捻りすると
$\nu_*(\omega_S(2H)\otimes\pi^*L_B^{-1})\hookrightarrow M\hookrightarrow\nu_*\pi^*L_B$。
両端は正規化線束で、合成は分数体上の乗法なので正規化に関して線形である。
$\omega_S(2H)=\pi^*(\omega_B\det E)$ から乗法切断は底に降り、$\deg E=8$ より次数は $2a$。
中央の切断が四つの核座標で尽きることは二次表示から従う。
Riemann--Rochによりこの包含が与える $U$ の独立切断数は少なくとも $4-a$ 本である。
特に $a=1$ なら長さ二の $\Theta$ での $U$ の評価階数は高々一で、核曲線への射は閉埋込みではない。
$a=2$ なら長さ四の因子が核内の直線上に写る（評価階数は二以下）。

核は非退化なので $e\ge3$、三次なら有理正規三次で二項四次Bezout型となり
定理\ref{thm:unrestricted-bezout-exclusion}に反する。
非退化な整四次空間曲線の算術種数は一以下である。種数一なら $h^0(\mathcal O_K(2))=8$ で、
核曲線は二つの独立な二次式に入る。これは六次余因子に二独立な定数関係を与える。
実際 $D$ 上では核の平方で、整八次上で消える六次式は恒等零である。
整八次の余因子が九または十独立であることに反する。算術種数零なら滑らかな有理四次、
$a=1$ であり、先の長さ二の評価は閉埋込みに反する。よって曲線核で $HL\le4$ はない。
\end{proof}

\begin{proof}[曲面核と一般平面]
核像は定数ではない（一般可逆な行列の余因子の定数核はない）。
曲面核なら $H^2=8$、$L^2$ は非退化な核曲面の次数の正の倍数で二以上である。
Hodge指数により $HL\le4$ は $HL=4$, $L^2=2$, $H\equiv2L$ しか許さない。
核射は二次曲面 $Q$ に双有理で、二つの射が潰す曲線は一致する。
定理\ref{thm:quadric-kernel-projection}の正規化とHilbert多項式の証明はこの数値等号だけで成立する。
すなわち二つの連結繊維の縮約を比較して $D$ の正規化は $Q$ となり、
入力束は $\mathcal O_Q(2)$、評価による包含の両辺のEuler多項式は $4(t+1)^2$ なので
$M=\nu_*\mathcal O_Q(1)$。これは定理\ref{thm:quintic-web-curve-excess}に反する。

一般入力平面切断の正規化上で余因子は六次式による核平方座標である。
従って共通固定部の全次数は $48-2HL\le38$。
各横断点で全余因子は極大イデアルの $r_i$ 乗に入るため、全正規化枝でのその固定次数は
少なくとも $r_i m_i$ である。一般一次式の枝位数の和が平面曲線の重複度になるためで、
節点や一般被約性は要らない。二行単元を消去できる一般階数二では $r_i\ge1,m_i\ge2$。
一般階数一では一行消去後の三行が全て極大イデアルに入り、$r_i\ge2,m_i\ge3$。
これらを合計して結論を得る。
\end{proof}
\path{computations/verify_sectional_cofactor_budget.py} は跡の次数、数値等号と横断次数の境界を検算する。
